\chapter{Fixed Income and RFQ Platform Connectivity}\label{ch:nw:fixed-income-and-rfq-platform-connectivity}

On a request-for-quote platform the client's \omterm{def:nw:fixed-income-and-rfq-platform-connectivity:timer}{quote timer} runs for seconds, not microseconds; yet dealers answer in
milliseconds. A response that arrives after the timer is not a worse price but no price at all, and some clients no longer
wait for the timer: their rules execute the best of the first few answers, or the first acceptable one. Under those rules a
dealer twenty milliseconds slower than a rival loses trades it would have won on price, and the chapter's model finds that
halving an auto-responder's pricing time is worth nearly fifteen times more than bidding a cent better.

Book~2 described request-for-quote trading, dealer-to-client platforms, inter-dealer brokers and all-to-all trading;
Book~11 built the auto-quoter with its win-probability model and cover price; Book~13 the FIX session and the latency budget.
This chapter connects them: the platforms' interfaces, the auto-responder's pipeline from request to answer, the latency that
matters when the clock is a \omterm{def:nw:fixed-income-and-rfq-platform-connectivity:timer}{quote timer}, and the inter-dealer order books, which are exchanges in all but name and sit in the
data centres of chapter~10 and~11. The auction is a labelled simulation on Book~11's price model.

\section{Platform interfaces}

A dealer connects to each dealer-to-client platform through a FIX session or the platform's own interface, receives requests,
and answers them with firm prices valid for a stated time. The platforms' protocols differ in what the client sees and how long
the dealer's price binds (\cref{dat:nw:fixed-income-and-rfq-platform-connectivity:platforms}): a one-sided request, a two-sided
request for market that hides the client's side, a request for stream in which dealers update their prices every 250
milliseconds, or a price with a ``wire time'' during which it can be hit.

\begin{definition}[Quote timer]\label{def:nw:fixed-income-and-rfq-platform-connectivity:timer}
A \emph{quote timer}\index{quote timer} is the period, set by the platform or the client for each request for quote, during
which dealers' answers are collected and remain executable; an answer that arrives after it expires is not shown, and a price
shown during it binds the dealer until the client trades, passes or the timer runs out.
\end{definition}

\begin{dated}{2026-09}{Request-for-quote protocols and automation, from the platforms}\label{dat:nw:fixed-income-and-rfq-platform-connectivity:platforms}
\begin{itemize}
\item TW SEF's protocols (advisory of 2015): an RFQ is ``a real-time auction with multiple RFQ recipients'' in which the
  requester selects the best price; at least three unaffiliated recipients for required transactions; request for market; RFQ+
  beside the order book; request for stream, with prices updated ``every 250 milliseconds''; otherwise a price with a ``wire
  time''.
\item Tradeweb's AiEX executes orders from the client's order management system by pre-defined rules; a client might ``execute at
  the best price once they get five responses'' or wait for a particular dealer (2018).
\item MarketAxess (June 2023): automation protocols were ``nearly 10\% of total trading volume'' in the year to date; its Adaptive
  Auto-X was in pilot.
\end{itemize}
\end{dated}

\section{Auto-responding: the pipeline from request to answer}

An auto-responder is a pipeline: the request arrives over the network and the FIX engine; the pricer values the bond (Book~11's
fair value from stale prints, the issuer curve and comparables); limits and risk are checked; the answer goes back. Each stage
has a latency distribution, and requests the auto-responder is not allowed to answer, for size, liquidity or risk, go to a
person, who answers in seconds.

\omcode{code/firm/rfqlink/firm_rfqlink.py}{47}{53}{An answer's time: the stages' lognormal times, plus a person's for the share outside the auto-responder's limits.}\label{lst:nw:fixed-income-and-rfq-platform-connectivity:response}

The chapter's dealer has five stages (network in and out \qty{0.5}{\milli\second} each, FIX engine \qty{0.1}{\milli\second},
pricing \qty{20}{\milli\second}, limits \qty{1}{\milli\second}, medians with lognormal spreads) and sends 2\% of requests to a
person, whose answers take six seconds at the median. Its four competitors are identical except for a pricing stage of
\qty{50}{\milli\second}. All of it is assumption, in the range an auto-responder can be built to (Book~13).

\begin{definition}[Automated execution rule]\label{def:nw:fixed-income-and-rfq-platform-connectivity:rule}
An \emph{automated execution rule}\index{automated execution rule} is a client's pre-set instruction, run by the platform, that
executes a request for quote without a person's decision: at the best answer when the timer expires, at the best of the first
$k$ answers, at the first answer within a tolerance of a reference price, or with a named dealer.
\end{definition}

\section{Latency that matters when the clock is a quote timer}

Two latencies matter. The first is whether the answer arrives before the timer: with a two-second timer the automatic stages never
miss, and the chance of missing is almost exactly the share sent to a person times the chance a person takes more than two seconds,
1.94\% for the chapter's dealer (\cref{fig:nw:fixed-income-and-rfq-platform-connectivity:miss}). The second is whether the answer
arrives before the answers that the client's rule will take, which depends on the rule.

\begin{omfigure}
\begin{tikzpicture}
\begin{axis}[width=11cm,height=5.6cm,xmode=log,ymode=log,xmin=100,xmax=31623,ymin=1e-5,ymax=0.1,
  xlabel={quote timer (ms, log)},ylabel={chance of missing it (log)},tick label style={font=\footnotesize},
  label style={font=\small},grid=major,grid style={gray!25},legend style={font=\footnotesize,at={(0.5,-0.3)},anchor=north,legend columns=2},
  unbounded coords=discard,table/col sep=comma]
\addplot[omThm,thick,mark=*,mark size=1.3pt] table[x=timer_ms,y=ours]{figdata/networks/25-fixed-income-and-rfq-platform-connectivity/miss.csv};
\addplot[omDef,thick,mark=square*,mark size=1.3pt] table[x=timer_ms,y=auto]{figdata/networks/25-fixed-income-and-rfq-platform-connectivity/miss.csv};
\legend{2\% answered by a person,fully automatic}
\end{axis}
\end{tikzpicture}

\omcaption{Simulation: the chance that the chapter's dealer answers after the \omterm{def:nw:fixed-income-and-rfq-platform-connectivity:timer}{quote timer}, against the timer's length. Fully
automatic, it never misses a timer of one second or more; with 2\% of requests answered by a person (six seconds at the median), it
misses 1.94\% of two-second timers, and the plateau is the person, not the network. Data: \texttt{fig\_rfq.py},
\texttt{nw\_rfq.miss\_curves()}.}\label{fig:nw:fixed-income-and-rfq-platform-connectivity:miss}
\end{omfigure}

The client's rule decides the rest. The auction draws each dealer's bid as Book~11 does, the dealer's noisy estimate of the bond's
value less a markup of 15 cents per 100 face, and the client sells only at or above its reservation price; with every answer in time
and the rule ``best at expiry'', the simulation is Book~11's auction exactly (the tests check it). Then the rules part ways.

\omcode{code/firm/rfqlink/firm_rfqlink.py}{68}{85}{The three client rules, and the requests lost to speed: our answer was the best acceptable one in time, and the rule chose another.}\label{lst:nw:fixed-income-and-rfq-platform-connectivity:rules}

\begin{proposition}[When speed is worth nothing]\label{prop:nw:fixed-income-and-rfq-platform-connectivity:speed}
Under the rule ``best at expiry'', a dealer's win rate depends on its response time only through the chance of missing the timer;
under ``best of the first $k$'' and ``first acceptable'', it falls with the dealer's rank in arrival order, so a faster answer at the
same price wins more.
\end{proposition}

\begin{proof}
At expiry the client compares all answers that arrived in time, whatever their order. Under the other rules, an answer outside the
first $k$, or behind an earlier acceptable one, is never compared, and a faster pipeline raises its chance of being among them.
\end{proof}

\begin{omfigure}
\begin{tikzpicture}
\begin{axis}[width=11cm,height=6cm,xmode=log,xmin=1,xmax=1000,ymin=0,ymax=0.65,
  xlabel={our pricing stage's median (ms, log); competitors at 50 ms},ylabel={share of requests we win},
  tick label style={font=\footnotesize},label style={font=\small},grid=major,grid style={gray!25},
  legend style={font=\footnotesize,at={(0.5,-0.3)},anchor=north,legend columns=3},table/col sep=comma]
\addplot[omDef,thick,mark=*,mark size=1.2pt] table[x=pricing_ms,y=expiry]{figdata/networks/25-fixed-income-and-rfq-platform-connectivity/wins.csv};
\addplot[omMeth,thick,mark=square*,mark size=1.2pt] table[x=pricing_ms,y=first_k]{figdata/networks/25-fixed-income-and-rfq-platform-connectivity/wins.csv};
\addplot[omThm,thick,mark=triangle*,mark size=1.6pt] table[x=pricing_ms,y=first_ok]{figdata/networks/25-fixed-income-and-rfq-platform-connectivity/wins.csv};
\legend{best at expiry,best of first three,first acceptable}
\end{axis}
\end{tikzpicture}

\omcaption{Simulation: the share of requests the dealer wins among five, against its pricing time, under three client rules, all
dealers at the same markup. At expiry the curve is flat until the pricing time's tail reaches the timer, at hundreds of milliseconds; under the other two rules every
millisecond counts. Data: \texttt{fig\_rfq.py}, \texttt{nw\_rfq.curves()}.}\label{fig:nw:fixed-income-and-rfq-platform-connectivity:wins}
\end{omfigure}

At its \qty{20}{\milli\second} pricing stage, the dealer wins 19.2\% of requests when clients take the best answer at expiry, 26.4\%
when they take the best of the first three, and 39.0\% when they take the first acceptable answer: being faster than its four rivals is
worth twice its share. The other side of the same rules is the business lost to speed: under ``first acceptable'', 5.7\% of all requests
are ones where the dealer's answer was the best acceptable one in time and a faster rival's was taken; under ``best of the first three'',
1.6\%.

\begin{table}[ht]
\centering\small
\begin{tabular}{@{}lrrr@{}}
\toprule
Client rule & Win rate & Halving pricing to 10 ms & One cent better \\
\midrule
Best at expiry & 19.2\% & $+0.0$ points & $+0.9$ points \\
Best of first three & 26.4\% & $+2.3$ points & $+1.0$ points \\
First acceptable & 39.0\% & $+12.6$ points & $+0.9$ points \\
\bottomrule
\end{tabular}
\caption{Simulation: what the dealer gains in win rate from halving its pricing time, against bidding one cent per 100 face (one basis
point of price) better, under each client rule. Data: \texttt{nw\_rfq.values()}.}\label{tab:nw:fixed-income-and-rfq-platform-connectivity:value}
\end{table}

The value of a millisecond and the value of a basis point are therefore not comparable in general but only per rule
(\cref{tab:nw:fixed-income-and-rfq-platform-connectivity:value}). A dealer whose clients wait for the timer should spend on price; one
whose clients execute the first acceptable answer should spend on its pipeline first, since ten milliseconds are worth nearly fifteen cents of
price. Which rules its clients use is information the dealer can estimate from its own wins and losses by response time.

\section{The interdealer order books}

The interdealer markets in government bonds are central limit order books, run like exchanges and located in the same campuses as
the futures and options markets of chapters~10 and~11 (\cref{dat:nw:fixed-income-and-rfq-platform-connectivity:idb}). BrokerTec's
U.S. markets match in CME's NY5 in Secaucus and its European markets in LD4.2 in Slough, each with a recovery site; BrokerTec Chicago, a
second order book for the actively traded Treasuries, sits in CME's Aurora data centre next to the Treasury futures, for trading cash
against futures. On these books everything in chapters~1 to~9 applies: \omterm{def:nw:colocation-products-and-how-they-are-sold:colo}{colocation}, \omterm{def:nw:colocation-products-and-how-they-are-sold:mmr}{cross-connects}, feeds, gateways; the RFQ platforms'
seconds become microseconds again.

\begin{dated}{2026-09}{Inter-dealer government bond order books on CME Globex, from CME's documentation}\label{dat:nw:fixed-income-and-rfq-platform-connectivity:idb}
BrokerTec US markets: production in NY5, recovery in Aurora; EU markets: production in LD4.2, recovery in NY5; recovery time objectives
of four hours (US) and two hours (EU). BrokerTec Chicago (secondary U.S. Treasury actives order book): the CME Globex data centre in
Aurora, Illinois.
\end{dated}

The distance between NY5 and Aurora is the cash--futures basis trader's distance: the U.S. Treasury futures match in Aurora, the
primary cash book in Secaucus, and the two are chapter~10's New Jersey--Chicago corridor apart, the route on which fibre and microwave
race. A second cash book in Aurora puts both legs in one building for the firms that trade them together.

\begin{method}[Connecting a fixed income desk]\label{met:nw:fixed-income-and-rfq-platform-connectivity:desk}
\begin{enumerate}
\item List the platforms and books the desk trades; for each, the protocol, the session type, the timer or wire time, and the data
  centre.
\item Build the auto-responder as a pipeline with a latency budget per stage; measure the share sent to a person and that person's
  times.
\item Estimate the clients' execution rules from wins and losses by response time; spend on speed where rules reward it and on price
  where they do not.
\item Colocate for the order books as for any exchange; for cash against futures, choose between two sites and one.
\item Monitor timer misses, answers after the first $k$, and wins lost to speed, per platform and per client.
\end{enumerate}
\end{method}

\section{Tutorial: two seconds to answer}\label{tut:nw:fixed-income-and-rfq-platform-connectivity:tut}

\begin{tutorial}
\textbf{Goal.} Simulate an auto-responder's pipeline, the timer, and three client rules on Book~11's auction.
\textbf{End state:} \cref{fig:nw:fixed-income-and-rfq-platform-connectivity:wins} and
\cref{tab:nw:fixed-income-and-rfq-platform-connectivity:value}.
\begin{tutsteps}
\item \textbf{Pipelines.} \texttt{firm\_rfqlink.Pipeline} and \texttt{Stage}; \texttt{response\_ms}
  (\cref{lst:nw:fixed-income-and-rfq-platform-connectivity:response}) and \texttt{miss\_prob} for the timer.
\item \textbf{Auction.} \texttt{auction} with each rule (\cref{lst:nw:fixed-income-and-rfq-platform-connectivity:rules}), checked
  against \texttt{firm\_rfqmm.simulate} when every answer is in time.
\item \textbf{Value.} \texttt{value\_of} for a faster pipeline and a better price; \texttt{nw\_rfq.curves} across pricing times.
\end{tutsteps}
\textbf{What to change next.} Let each client choose its rule at random with shares estimated from data, and let the markup respond:
Book~11's equilibrium under a first-acceptable rule.
\end{tutorial}

\section{Build: the RFQ connectivity model}\label{bld:nw:fixed-income-and-rfq-platform-connectivity:rfqlink}

\begin{build}
\bfield{Purpose} Timers, pipelines and client rules for a dealer on RFQ platforms: the fixed income rows of chapter~29's plan.
\bfield{Interface} \texttt{firm\_rfqlink}: \texttt{Stage}, \texttt{Pipeline}, \texttt{response\_ms}, \texttt{miss\_prob},
\texttt{auction}, \texttt{value\_of}, \texttt{scaled}, \texttt{log\_grid}.
\bfield{Rules} Prices follow Book~11's \texttt{firm.rfqmm} model, unchanged; pipeline and client parameters are stated assumptions;
results are a labelled simulation.
\bfield{Acceptance tests} \texttt{code/firm/rfqlink/tests/}: identity with \texttt{firm\_rfqmm.simulate} when every answer is in time,
deterministic pipelines and timers, and the rules' ordering.
\bfield{Stretch} Wire times and last look on streams; clients' rule mix estimated from data; throttles on sessions.
\end{build}

\begin{omsources}
\item TW SEF LLC, Trading and Execution Protocols (2015).
\item The TRADE, ``Automating trade execution, intelligently'' (2018); Markets Media on MarketAxess's Adaptive Auto-X (2023).
\item CME Group Client Systems Wiki: BrokerTec disaster recovery process and scenarios; BrokerTec Chicago.
\end{omsources}

\section{Exercises}

\begin{exercise}[$\star$]\label{exo:nw:fixed-income-and-rfq-platform-connectivity:1}
What happens to an answer that arrives after the \omterm{def:nw:fixed-income-and-rfq-platform-connectivity:timer}{quote timer}?
\end{exercise}

\begin{exercise}[$\star$]\label{exo:nw:fixed-income-and-rfq-platform-connectivity:2}
Where do BrokerTec's U.S. and European order books match, and where do they recover?
\end{exercise}

\begin{exercise}[$\star$]\label{exo:nw:fixed-income-and-rfq-platform-connectivity:3}
Why is the chance of missing a two-second timer almost the share of requests sent to a person?
\end{exercise}

\begin{exercise}[$\star\star$]\label{exo:nw:fixed-income-and-rfq-platform-connectivity:4}
Using \cref{prop:nw:fixed-income-and-rfq-platform-connectivity:speed}, why does halving the pricing time gain nothing under ``best at
expiry''?
\end{exercise}

\begin{exercise}[$\star\star$]\label{exo:nw:fixed-income-and-rfq-platform-connectivity:5}
A client executes ``at the best price once it has five responses'' and asks five dealers. Which of the chapter's rules is that?
\end{exercise}

\begin{exercise}[$\star\star$]\label{exo:nw:fixed-income-and-rfq-platform-connectivity:6}
Why might a cash--futures basis trader prefer a second Treasury order book in Aurora?
\end{exercise}

\begin{exercise}[$\star\star\star$]\label{exo:nw:fixed-income-and-rfq-platform-connectivity:7}
\emph{Coding.} With \texttt{nw\_rfq.curves}, at what pricing time does the dealer's win rate under ``first acceptable'' fall below its
rate under ``best at expiry''?
\end{exercise}

\begin{exercise}[$\star\star\star$]\label{exo:nw:fixed-income-and-rfq-platform-connectivity:8}
\emph{Find the flaw.} ``Our clients have two-second timers, so our 20-millisecond pricer is a hundred times faster than it needs to be.''
\end{exercise}

\section{Problem: Two Seconds to Answer}

\begin{problem}[{Weekend problem --- timers, rules and a dealer's pipeline}]\label{pb:nw:fixed-income-and-rfq-platform-connectivity:1}
A dealer answers requests for quote on corporate bonds against four competitors. Its pipeline and theirs are the chapter's; the timer
is two seconds.

\medskip\noindent\textbf{Part I --- The timer.}
\begin{enumerate}
\item What is the chance that the dealer misses the timer, and where does it come from?
\item What would it be fully automatic?
\item At what timer length does the fully automatic pipeline stop missing?
\item What does a missed answer cost the dealer, compared with a losing one?
\end{enumerate}

\medskip\noindent\textbf{Part II --- The rules.}
\begin{enumerate}[resume]
\item What share of requests does the dealer win under each rule?
\item What share is lost to speed under each?
\item Why is it zero at expiry?
\item Why does the dealer win more under ``first acceptable'' than at expiry?
\end{enumerate}

\medskip\noindent\textbf{Part III --- Speed against price.}
\begin{enumerate}[resume]
\item What does halving the pricing time gain under each rule?
\item What does bidding one cent better gain?
\item How many cents of price are ten milliseconds worth under ``first acceptable''?
\item How would the dealer find out which rules its clients use?
\end{enumerate}

\medskip\noindent\textbf{Part IV --- The verdict.}
\begin{enumerate}[resume]
\item State the \emph{named result}: the probability of missing the \omterm{def:nw:fixed-income-and-rfq-platform-connectivity:timer}{quote timer} at the pipeline's latency distribution, and the share
  of business lost to a faster dealer under a first-acceptable execution rule.
\item Which inputs are published and which assumed?
\item Where should the dealer spend its next engineering month?
\item What happens if every competitor also halves its pricing time?
\item How should the share sent to a person be managed?
\item What changes on the inter-dealer order books?
\item Where does this go in chapter~29's plan?
\item In one sentence: when does a millisecond matter on a platform whose timer is two seconds?
\end{enumerate}
\end{problem}

\section{Interview questions}

\begin{interviewq}[$\star$ \iqroles{developer}]\label{iq:nw:fixed-income-and-rfq-platform-connectivity:1}
What are the stages of an RFQ auto-responder, and which one would you expect to be slowest?
\end{interviewq}

\begin{interviewq}[$\star\star$ \iqroles{trader}]\label{iq:nw:fixed-income-and-rfq-platform-connectivity:2}
Why would a client execute the first acceptable answer rather than wait for the best?
\end{interviewq}

\begin{interviewq}[$\star\star$ \iqroles{developer, researcher}]\label{iq:nw:fixed-income-and-rfq-platform-connectivity:3}
How would you estimate whether your response time costs you business on a platform?
\end{interviewq}

\begin{interviewq}[$\star\star$ \iqroles{developer}]\label{iq:nw:fixed-income-and-rfq-platform-connectivity:4}
Your auto-responder's p99 latency doubled after a release. What do you check?
\end{interviewq}

\begin{interviewq}[$\star\star$ \iqroles{researcher}]\label{iq:nw:fixed-income-and-rfq-platform-connectivity:5}
Should a dealer quote tighter or answer faster? How would you decide?
\end{interviewq}

\begin{interviewq}[$\star\star\star$ \iqroles{developer, researcher}]\label{iq:nw:fixed-income-and-rfq-platform-connectivity:6}
Design the connectivity and auto-responder for a dealer answering on three bond platforms and trading on an inter-dealer book.
\end{interviewq}
